% =====================================================================================
%  Data, label construction and the admissibility protocol.
%  Every number here comes from code/results_deep_7030/dataset_stats.csv and
%  priority_breakdown.csv, written by evaluation/deep_metrics.py.
% =====================================================================================

\section{Data and the admissibility protocol}
\label{sec:data}

\subsection{The log}

BPI Challenge 2014 \citep{dongen2014challenge} is the incident-management record of Rabobank
Nederland Group ICT, released for the 2014 Business Process Intelligence Challenge. It is one
of the few public logs in which the assignment decision, the resource that received it and the
eventual outcome are all recorded on real work rather than simulated. Three files are used: an
incident header, an activity log and an interaction log. Table~\ref{tab:dataset} reports the
working statistics.

\begin{table}[htbp]
\centering
\caption{Dataset statistics. Counts are produced by \texttt{evaluation/deep\_metrics.py} and
written to \texttt{dataset\_stats.csv}; nothing in this table is typed by hand.}
\label{tab:dataset}
\begin{tabular}{lr}
\toprule
Property & Value \\
\midrule
Raw incident records & \num{31238} \\
Raw activity events & \num{210837} \\
Raw interaction records & \num{68358} \\
Incidents with a resolvable handle time and priority & \num{31236} \\
Incidents in the predictive modelling matrix & \num{31236} \\
Distinct first-assignment groups & \num{117} \\
Distinct configuration-item types & \num{13} \\
Distinct configuration-item subtypes & \num{58} \\
Distinct service components (WBS codes) & \num{259} \\
Admissible attributes after the protocol & \num{17} \\
Attributes excluded as inadmissible & \num{9} \\
\midrule
Breach cases (positive class) & \num{11444} \\
Non-breach cases (negative class) & \num{19792} \\
Positive class rate & \num{0.3664} \\
\midrule
Training incidents, random \num{70}/\num{30} split & \num{21865} \\
Test incidents, random \num{70}/\num{30} split & \num{9371} \\
Training incidents, chronological holdout & \num{24988} \\
Test incidents, chronological holdout & \num{6248} \\
\midrule
Earliest incident open time & 2012-01-10 \\
Latest incident open time & 2014-12-01 \\
Observation window & \num{1055} days \\
\bottomrule
\end{tabular}
\end{table}

Four properties of the raw files silently corrupt a naive load, and they are recorded here
because they cost real debugging time and because a reader attempting to reproduce this work
will meet all four. The files are semicolon delimited rather than comma delimited. They are
encoded latin-1 rather than UTF-8, so the group names carry mojibake under a UTF-8 read. The
handle-time column uses a decimal comma, so a direct numeric cast silently produces nulls.
And the incident header pads to 78 columns with empty trailing delimiters, which shifts the
column index of every field a naive parser reads positionally.

\subsection{Two working sets, and why they differ}
\label{sec:two-sets}

The paper has two components with different data requirements, and they operate on different
subsets of the same log. Reporting one figure for both would be wrong, so both are stated.

The \emph{predictive} component uses every incident with a resolvable handle time and a
recorded priority, which is \num{31236} incidents across \num{117} distinct first-assignment
groups. It places no constraint on group size, because a target encoding with a smoothing prior
handles a rarely used group correctly by shrinking it towards the population rate.

The \emph{assignment} component of Section~\ref{sec:assignment-results} simulates a decision
over a candidate set, and a candidate set containing groups that have handled three incidents
in three years is not a realistic decision. It therefore restricts to assignment groups with at
least \num{50} recorded cases, which leaves \num{15828} incidents, \num{50} groups and
\num{47} configuration-item categories, of which the earliest \num{12662} by open time form the
training period on which every parameter is fitted.

\subsection{The service level label}
\label{sec:label}

The log's own \texttt{Alert Status} field reads \texttt{closed} on \num{46606} of the
\num{46809} raw records and carries no usable breach signal, so no contractual service level
target is recoverable from the data. The label is therefore constructed, and the construction
is stated precisely rather than described loosely.

Let $H_i$ be the recorded handle time of incident $i$ in hours and $\pi(i)$ its priority class.
Write $\tilde{H}_{p}$ for the median handle time of priority class $p$. An incident breaches
when
\begin{equation}
y_i \;=\; \mathbb{1}\bigl[\, H_i \;>\; \tau_{\pi(i)} \,\bigr],
\qquad \tau_{p} \;=\; 2\,\tilde{H}_{p}.
\label{eq:label}
\end{equation}

The factor of two is a deliberate choice of a lenient target: at a factor of one the positive
class would be half the data by construction and the label would carry no information beyond
the median. Table~\ref{tab:priority} gives the resulting thresholds. The breach rate is
notably flat across priority classes, between \SI{33.6}{\percent} and \SI{37.6}{\percent},
which is a direct consequence of defining the threshold within each class: the label measures
whether a case was slow \emph{for its own priority}, not whether it was slow in absolute terms.

\begin{table}[htbp]
\centering
\caption{Per-priority volume, median handle time, derived service level threshold and breach
rate. Priority 1 does not occur in the filtered log.}
\label{tab:priority}
\begin{tabular}{lrrrr}
\toprule
Priority & Incidents & Median handle (h) & Threshold $\tau_p$ (h) & Breach rate \\
\midrule
2 & \num{521}   & 3.68 & 7.36  & 0.3359 \\
3 & \num{4347}  & 1.58 & 3.17  & 0.3759 \\
4 & \num{14810} & 4.05 & 8.10  & 0.3629 \\
5 & \num{11558} & 8.37 & 16.74 & 0.3686 \\
\bottomrule
\end{tabular}
\end{table}

\begin{figure}[htbp]
\centering
\includegraphics[width=0.94\textwidth]{fig10_class_balance.pdf}
\caption{Class balance and priority composition. The positive class is a
\SI{36.6}{\percent} minority, which is imbalanced enough that accuracy alone would be
misleading and is why precision, recall, the Matthews correlation coefficient and average
precision are reported throughout.}
\label{fig:classbalance}
\end{figure}

\paragraph{Where the threshold is estimated, and where that matters.}
For the random-split arms, $\tilde{H}_p$ is computed over the whole log. This is a property of
the \emph{label}, not of a feature: no model observes $\tilde{H}_p$, and every model is scored
against the same fixed label, so the comparison between configurations is unaffected. For the
chronological arm it would nevertheless be wrong, because a system deployed in January 2014
cannot know the 2014 medians. That arm therefore re-derives $\tau_p$ from training-period cases
alone, and the effect is reported in Section~\ref{sec:temporal}.

\subsection{The protocol}
\label{sec:protocol}

Let $t^{\mathrm{dec}}_i$ be the moment the assignment decision for case $i$ is taken, and let
$\mathcal{F}_{t}$ be the information available to the organisation at time $t$.

\paragraph{Rule 1, feature admissibility.} An attribute $x^{(k)}$ is admissible for case $i$
only if its value is $\mathcal{F}_{t^{\mathrm{dec}}_i}$-measurable:
\begin{equation}
x^{(k)}_i \in \mathcal{F}_{t^{\mathrm{dec}}_i}
\quad\text{for every case } i .
\label{eq:rule1}
\end{equation}
Any attribute aggregated over the whole closed case violates this by construction, whatever its
apparent predictive value. This is \citeauthor{kaufman2012leakage}'s legitimacy condition
\citep{kaufman2012leakage} instantiated at a specific decision point rather than stated in
general.

\paragraph{Rule 2, encoding admissibility.} A group-level target encoding for group $g$
evaluated for case $i$ may use only cases resolved strictly before case $i$ opened:
\begin{equation}
\hat{\theta}_g(i) \;=\;
\frac{\sum_{j \in \mathcal{P}(i,g)} y_j \;+\; \alpha \bar{y}_{\mathcal{P}}}
     {\lvert \mathcal{P}(i,g) \rvert \;+\; \alpha},
\label{eq:rule2}
\end{equation}
where the admissible past set is
\begin{equation}
\mathcal{P}(i,g) = \bigl\{\, j : g_j = g,\; t^{\mathrm{res}}_j < t^{\mathrm{open}}_i \,\bigr\},
\label{eq:pastset}
\end{equation}
with smoothing prior $\alpha = 20$ and $\bar{y}_{\mathcal{P}}$ the prior breach rate over the
same past set. Operationally this means the encoding is refitted inside every cross-validation
fold and again on the outer training split, never once over the full data before splitting.
\citet{pargent2022regularized} show that regularised target encoding outperforms conventional
categorical handling; the regularisation is exactly what makes refitting necessary, because an
encoding fitted once over everything contains each row's own label inside its own predictor.

Applying both rules to BPI 2014 admits \num{17} attributes and excludes nine.
Table~\ref{tab:attributes} lists every one with the reason for its verdict. The table is
generated from \texttt{feature\_\allowbreak justifications\_\allowbreak leakfree.csv}, which the pipeline writes with a
\texttt{Known\_\allowbreak at\_\allowbreak assignment\_\allowbreak time} column, so the verdict is a machine-checked property of
the code rather than a claim in prose.

\begin{table}[htbp]
\centering
\caption{Attribute admissibility. Seventeen attributes are admitted, nine excluded. The two
target encodings are admissible only under Equation~\ref{eq:rule2}; used the way version one of
this pipeline used them, they are the single largest source of contamination.}
\label{tab:attributes}
\small
\begin{tabular}{L{4.0cm} c L{8.4cm}}
\toprule
Attribute & Admitted & Reason \\
\midrule
Priority & yes & Impact times urgency as recorded at open. The service level clock is set per priority band. \\
Impact & yes & Business impact level recorded at open. \\
Urgency & yes & Urgency level recorded at open, independent of impact. \\
Open hour & yes & Hour of arrival. Proxies shift staffing, which governs pick-up speed. \\
Open day of week & yes & Weekday of arrival. Captures the weekly demand cycle. \\
Is weekend & yes & Reduced staffing materially changes achievable response time. \\
Is business hours & yes & Out-of-hours arrivals queue until staffed. \\
Open month & yes & Seasonal load and release-cycle effects. \\
Queue length at open & yes & Incidents already open at arrival. Computed from other cases only. \\
Assignment delay hours & yes & Elapsed hours from open to first assignment. Observable at the decision; not total handle time. \\
CI type & yes & Asset class of the affected configuration item. \\
CI subtype & yes & Finer-grained asset class. \\
Service component (WBS) & yes & Work breakdown code of the affected service. \\
Category & yes & Incident category assigned at intake. \\
Alert status & yes & Alert state on the record at open. \\
Group breach rate (encoded) & yes$^{\dagger}$ & Smoothed historical breach rate of the first assignment group, past cases only, prior 20. \\
Group volume (encoded) & yes$^{\dagger}$ & Case volume of the group, past cases only. Proxies capacity and experience. \\
\midrule
Number of reassignments & \textbf{no} & Total over the whole closed case. Rises with duration, which defines the label. \\
Groups touched & \textbf{no} & Distinct groups over the whole closed case. Post hoc. \\
Total activity events & \textbf{no} & Every activity event across the closed case. Strongest duration proxy present. \\
Has reopen & \textbf{no} & Reopens occur late in a case. Post hoc. \\
Related incidents & \textbf{no} & Final tally on the closed record, not the value visible at open. \\
Interaction count & \textbf{no} & Aggregated across the full case from the interaction log. Post hoc. \\
First-call-resolution rate & \textbf{no} & Resolution outcome is known only after handling. \\
Group mean handle time (v1 form) & \textbf{no} & Mean of the quantity the label thresholds, over the full data before the split. \\
Group breach rate (v1 form) & \textbf{no} & Mean of the target over the full data including test rows and the row's own label. \\
\bottomrule
\end{tabular}

\vspace{0.4em}
{\footnotesize $^{\dagger}$Admissible only under Equation~\ref{eq:rule2}. The same two
quantities appear in the excluded list in the form version one of this pipeline computed them.}
\end{table}
